% =============================================================================
% 文件: sections/s9_appendix.tex   附录
% =============================================================================
\newpage
\appendix
\renewcommand{\thesection}{附录\Alph{section}}

\section{支撑材料文件列表}

本论文的支撑材料为一个压缩包，其目录结构与文件清单如下
（所有源程序均为可直接运行的 Python 3 脚本，不含任何身份信息）。

\begin{center}
\small
\begin{longtable}{@{}p{4.5cm}p{9.8cm}@{}}
\toprule
文件 / 目录 & 内容说明 \\
\midrule
\endfirsthead
\toprule
文件 / 目录 & 内容说明 \\
\midrule
\endhead
\texttt{code/drying\_core.py} & 核心物理与数值模块：物性关系、环境与半径历史、有限体积求解器、表面重构、干燥终点判据 \\
\texttt{code/scenarios.py} & 四个问题的场景配置、结果 Excel 写入、位置换算 \\
\texttt{code/plotstyle.py} & 论文绘图统一风格（中文字体、分辨率） \\
\texttt{code/s01\_data\_analysis.py} & 第 2 步：附件数据探索与数据形态图 \\
\texttt{code/s02\_solve\_p1.py} & 问题 1 求解、表 1/表 2、网格与时间步收敛 \\
\texttt{code/s03\_solve\_p23.py} & 问题 2、3 求解、表 3/4/5、干燥时间与收敛证据 \\
\texttt{code/s04\_solve\_p4.py} & 问题 4 求解、表 6、收缩轨迹与半径插值对比 \\
\texttt{code/s05\_verification.py} & 解析解对照、独立方法互证、守恒检验、干物质守恒、合成数据反演 \\
\texttt{code/s06\_sensitivity.py} & 单因素扰动、平台口径、含湿量口径、噪声注入、Monte Carlo \\
\texttt{code/s07\_figures.py} & 论文全部插图的生成脚本 \\
\texttt{code/s08\_make\_tables.py} & 由结果数据自动生成论文中的 LaTeX 表格 \\
\texttt{data/input/} & 题目附件 1、附件 2 与结果模板（原始数据，未修改） \\
\texttt{results/result1.xlsx} & 问题 1 结果：\texttt{温度}、\texttt{水分浓度} 两个工作表，$1\ \mathrm{s}\times0.1\ \mathrm{cm}$ \\
\texttt{results/result2.xlsx} & 问题 2 结果：\texttt{温度}、\texttt{水分浓度} 两个工作表，$1\ \mathrm{s}\times0.1\ \mathrm{cm}$ \\
\texttt{results/result3.xlsx} & 问题 3 结果：$60\ \mathrm{s}\times0.1\ \mathrm{cm}$ \\
\texttt{results/result4.xlsx} & 问题 4 结果：$60\ \mathrm{s}\times0.1\ \mathrm{cm}$，末列为药材表面 \\
\texttt{data/*.csv} & 论文表格对应的机器可读结果、收敛性、灵敏度与反演数据 \\
\texttt{data/*.txt} & 各步骤的完整运行摘要（含守恒误差、收敛阶、检验结论） \\
\texttt{figures/} & 论文全部插图（PNG, 220 dpi） \\
\texttt{paper/} & 论文 LaTeX 源码、\texttt{cumcmthesis.cls} 模板与编译产物 \\
\texttt{docs/自检报告.md} & 逐步自检记录（运行结果、遇到的问题与解决方式） \\
\texttt{docs/参考文献来源.md} & 每篇参考文献的检索渠道、核实结论与支撑论点 \\
\bottomrule
\end{longtable}
\end{center}

\noindent 说明：压缩包内所有文件均不含任何显示参赛者身份、学校与赛区的信息。

% 源程序清单由 code/s09_package.py 自动生成, 保证不漏项、与支撑材料一致
% 由 s09_package.py 自动生成: 附录源程序清单, 请勿手工编辑
\section{完整源程序}

以下给出建模与求解所使用的\textbf{全部}源程序，内容与支撑材料中的.py 文件逐字节一致，可直接运行。
运行顺序：\texttt{s01} $\to$ \texttt{s02} $\to$ \texttt{s03} $\to$ \texttt{s04} $\to$ \texttt{s05} $\to$ \texttt{s06} $\to$ \texttt{s07}（另有 \texttt{s08}$\sim$\texttt{s17} 为出表、回填、编译、打包等辅助脚本）。

\subsection{核心求解模块 drying\_core.py}
\lstinputlisting[language=Python]{code_appendix/drying_core.txt}

\subsection{场景与结果输出 scenarios.py}
\lstinputlisting[language=Python]{code_appendix/scenarios.txt}

\subsection{绘图风格 plotstyle.py}
\lstinputlisting[language=Python]{code_appendix/plotstyle.txt}

\subsection{数据探索 s01\_data\_analysis.py}
\lstinputlisting[language=Python]{code_appendix/s01_data_analysis.txt}

\subsection{问题 1 求解 s02\_solve\_p1.py}
\lstinputlisting[language=Python]{code_appendix/s02_solve_p1.txt}

\subsection{问题 2、3 求解 s03\_solve\_p23.py}
\lstinputlisting[language=Python]{code_appendix/s03_solve_p23.txt}

\subsection{问题 4 求解 s04\_solve\_p4.py}
\lstinputlisting[language=Python]{code_appendix/s04_solve_p4.txt}

\subsection{模型验证 s05\_verification.py}
\lstinputlisting[language=Python]{code_appendix/s05_verification.txt}

\subsection{灵敏度与不确定度 s06\_sensitivity.py}
\lstinputlisting[language=Python]{code_appendix/s06_sensitivity.txt}

\subsection{论文插图 s07\_figures.py}
\lstinputlisting[language=Python]{code_appendix/s07_figures.txt}

\subsection{论文表格生成 s08\_make\_tables.py}
\lstinputlisting[language=Python]{code_appendix/s08_make_tables.txt}

\subsection{交付物打包 s09\_package.py}
\lstinputlisting[language=Python]{code_appendix/s09_package.txt}

\subsection{论文数字自动回填 s10\_fill\_numbers.py}
\lstinputlisting[language=Python]{code_appendix/s10_fill_numbers.txt}

\subsection{守恒检验复算 s11\_recheck\_conservation.py}
\lstinputlisting[language=Python]{code_appendix/s11_recheck_conservation.txt}

\subsection{结果文件标签补正 s12\_fix\_outputs.py}
\lstinputlisting[language=Python]{code_appendix/s12_fix_outputs.txt}

\subsection{论文编译与交付自检 s13\_build\_paper.py}
\lstinputlisting[language=Python]{code_appendix/s13_build_paper.txt}

\subsection{合成数据反演复算 s14\_inverse.py}
\lstinputlisting[language=Python]{code_appendix/s14_inverse.txt}

\subsection{反演稳健估计 s15\_inverse\_e1.py}
\lstinputlisting[language=Python]{code_appendix/s15_inverse_e1.txt}

\subsection{问题 4 插图 s16\_fig\_p4.py}
\lstinputlisting[language=Python]{code_appendix/s16_fig_p4.txt}

\subsection{正文页数压缩 s17\_compact\_paper.py}
\lstinputlisting[language=Python]{code_appendix/s17_compact_paper.txt}



\section{结果数据文件说明}

\texttt{results/} 下的四个 Excel 文件严格按附件 3 模板的字段约定填写：
每个工作表第 1 行为到药材中心的距离（单位 cm），第 1 列为时间（单位 s），
其余为对应物理量，全部保留四位小数。问题 \texttt{result1}/\texttt{result2}
包含 \texttt{温度} 与 \texttt{水分浓度} 两个工作表；
\texttt{result3}/\texttt{result4} 为单表，只保存水分浓度。
问题 4 因药材持续收缩，到药材中心的距离超过 $1.1\ \mathrm{cm}$ 后不再位于
药材内部（末态半径仅 $1.198\ \mathrm{cm}$），故末列统一记为“药材表面”，
其取值为 $r=R(t)$ 处的半单元重构值。
